
\subsection{Vehicular and Towed SAM Units}

\begin{figure*}
    \centering
    \includegraphics[width=0.8\linewidth]{figure-sam-launcher.png}
    \caption{SAM Launcher Units}
    \label{figure:sam-launcher-units}
\end{figure*}

\paragraph{Echelon}
The size of a SAM launcher unit is squad, section, or platoon according to its damage resilience (1D, 2D, or 3D, respectively).
This does not always coincide with the actual designations used historically. For example, the SA-2 battery in the game corresponds to a battalion in reality.

\paragraph{Nuclear Warheads}
The SA-2E has the option of a nuclear warhead.

\subsection{Soviet and Russian SAM Launcher Units}

\paragraph{SA-2B/C/E/F Battery}

This is an S-75 (SA-2 Guideline) battery of six single launchers for the long-range V-750 missile and their associated P-12 (Spoon Rest) EWR and SNR-75 (Fan Song) TTR. The variants are the S-75 Desna (SA-2B), S-75M Volkhov (SA-2C), S-75AK (SA-2E), and S-75SM (SA-2F).
All are radar-guided, but the SA-2F also has an optical backup.

The S-75 Dvina (SA-2A) system was deployed by the Soviet PVO in 1957. In addition to serving as a key part of the air-defense system of the Soviet Union (and its notable role in ending reconnaissance overflights by shooting down the U-2 of Gary Powers) and its allies, it saw combat in Cuba in 1962, with North Vietnamese forces from 1965, with Indian forces in the 1965 war with Pakistan, and with Egyptian and Syrian forces from 1967.

\paragraph{SA-3B Battery}

This is an S-125 Neva (SA-3B Goa) battery with four quad launchers for the 5V24 missile and their associated P-15 (Flat Face) EWR and SNR-125 (Low Blow) TTR.

It saw combat with Egyptian and Syrian forces from 1970.

\paragraph{Mobile SA-3B Battery}

This is a mobile S-125 Neva (SA-3B Goa) battery with four twin launchers for the 5V24 missile on trucks and their associated P-15 (Flat Face) EWR and SNR-125 (Low Blow) TTR.

It saw combat with Egyptian and Syrian forces from 1970 and with Iraqi forces in the Iran-Iraq War, the Gulf War, and the Invasion of Iraq.

\paragraph{SA-4A/B Battery}

This is an armored 2K11 Krug (SA-4A/B Ganef) battery with three launcher vehicles each with two long-range 9M8M1/9M8M2 missiles and their associated P-40 (Long Track) EWR and 1S32 (Pat Hand) TTR.

It saw combat with Egyptian and Syrian forces from 1970 and with Iraqi forces in the Iran-Iraq War, the Gulf War, and the Invasion of Iraq.

\paragraph{SA-6A/B Battery}

This is an armored 2K12 Kub (SA-6A/B Gainful) battery with four launcher vehicles, each with three medium-range 2K12 Kub (SA-6A) or 3M9M3 (SA-6B) missiles and their associated P-40 (Long Track) EWR and 1S32 (Pat Hand) TTR.

It saw combat with Egyptian and Syrian forces from 1970 and with Iraqi forces in the Iran-Iraq War, the Gulf War, and the Invasion of Iraq.

\paragraph{SA-8A/B Squad}

This is a squad with a single 9K33 Osa or 9K33M Osa-AK (SA-8A/B Gecko) vehicle with integrated 1S51M3 (Land Roll) TTR.
The SA-8A carries four 9M33 missiles, and the SA-8B carries six 9M33M missiles.

The Osa has been widely exported and seen combat with Syrian forces in the 1982 Lebanon War, with Cuban forces in the South African Border War, with Iraqi forces in the 1990 Gulf War, and in more recent conflicts.

\paragraph{SA-9A/B Squad}

This is a squad with a single 9K21 Strela-1 or 9K21M Strela-1M (SA-9A/B Gaskin) vehicle with four 9M21 or 9M21M missiles.
In the Soviet army, a platoon of four was assigned to the air-defense battery of tank and motor-rifle regiments, complementing the ZSU-23-4.

The Strela-1 has been widely exported and saw combat in the South African Border War, the Iran-Iraq War, the Lebanon War, the Gulf War, the Kosovo War, and the Invasion of Iraq.

\paragraph{SA-11 Squad}

This is a squad with a single 9K37 Buk (SA-11 Gadfly) vehicle with four 9M37 missiles and an integrated 9S35 (Fire Dome) TTR. Three squads make up a battery, and three batteries make up a battalion. Normally, one armored 9S18 (Tube Arm) or 9S18M1 (Snow Drift) EWR is allocated to each battalion.

The SA-11 has seen combat in the War in Abkhazia, the Russo-Ukrainian War, and during the Syrian Civil War.

\paragraph{SA-13 Squad}

This is a single 9K35 Strela-10 (SA-13 Gopher) vehicle with four infra-red homing 9M37 missiles. Four squads form a platoon.
It began to replace the SA-9B in the air-defense platoons of tank and motor-rifle regiments in the Soviet Army starting in 1976 and was subsequently widely exported to Soviet allies.
Compared to the earlier system, the missile offers greater lethality, and the vehicle greater protection and mobility (although it is not amphibious).
Some are equipped with 9S16 (Flat Box-B) RWR.

The SA-13 saw combat in the South African Border War with Angolan forces, the Gulf War with Iraqi forces, and in the Kosovo War with Serbian forces.

\subsection{US SAM Launcher Units}

\paragraph{M48 Squad}

This is a single M48 Chaparral vehicle with four adapted Sidewinder infrared-homing missiles. It was developed after the ambitious Mauler SAM program was abandoned and introduced in 1969. Night IR sights were added in 1984 to give it a limited night capability. The MIM-72A was the original missile based on the AIM-9D. The MIM-72C added an improved seeker and warhead, and the MIM-72E/F added a smokeless motor. The final MIM-72G added a seeker based on that of the FIM-92B Stinger, with much improved rejection of infrared countermeasures.

\paragraph{I-HAWK Platoon}

This consists of three triple launchers for the MIM-23A/C/D missile, plus supporting radars and command elements.
The MIM-23B Improved HAWK and its associated improved radars were developed from the original HAWK MIM-23A to give much improved performance at low levels.
The MIM-23C has improved ECCM, and the MIM-23D adds a home-on-jam capability.
Two platoons with an EWR make up a battery.

The I-HAWK was deployed by the US Army, US Marines, and many US allies.

\paragraph{MIM-104 Battery}

An MIM-104 Patriot battery consists of six launchers with associated radar and command elements.
The MIM-104A uses the “Standard” missile and was capable only of defense against aircraft.
Patriot is used by the US Army and numerous US allies.

\begin{table*}[p]
    \caption*{\hypertarget{table:sam-launcher-unit}{SAM Launcher Unit Table}}
    \centering
    \footnotesize
    \input{table-sam-launcher-units.tex}
    \label{table:sam-launcher-units}
\end{table*}
